\chapter{Conclusion and Future Work}

\section{Summary}

This project set out to translate American Sign Language gestures into speech
using only a wearable inertial glove, with the entire recognition pipeline
running on commodity hardware and a phone. The resulting system places a
microcontroller-and-IMU module on each fingernail, streams four fingers' data
over ESP-NOW to a fifth ``master'' finger, and bridges to a mobile phone over
Bluetooth Low Energy, where a quantised one-dimensional convolutional neural
network classifies a captured window and speaks the result. The wireless link
delivers all five finger streams at 25\,Hz with negligible loss, and the
deployed model runs in about one millisecond using tens of kilobytes of memory.

Beyond the working system, the project's most valuable outcome is a clear,
empirically grounded understanding of \emph{feature representation} for
inertial sign recognition. Through successive live failures and fixes, the work
established that absolute-orientation features for static handshapes are
fragile against imperfect re-calibration and limited by sensor resolution; that
an orientation-invariant relative transform fixes the calibration problem but
not the resolution problem, and harms dynamic signs that depend on global
motion; and that the robust formulation is to perform gestures as distinctive
\emph{motions} and classify them on \emph{absolute} features, placing the
decisive information in the calibration-independent angular-rate channels.

\section{Limitations}

Several limitations are acknowledged honestly. First, the reported accuracies
are within-session; rigorous cross-session and cross-user evaluation remains to
be done, and is expected to lower the figures. Second, the device instruments a
single hand and recognises segmented, one-gesture-at-a-time input, not
continuous or two-handed signing. Third, fingerspelling letters that differ
only by fine thumb placement are at the limit of what six-axis IMUs can resolve
statically; the project mitigates this by encoding letters as distinct motions,
but this departs from the canonical static handshapes. Fourth, power integrity
under simultaneous radio transmissions is a real constraint that requires
attention to the supply and wiring. Finally, recognition operates at the level
of single letters and words without higher-level language modelling.

\section{Future Work}

\begin{itemize}
    \item \textbf{Cross-session and multi-user data.} Collecting several
    sessions on different days, and data from multiple wearers, would both
    improve robustness and provide an honest generalisation measure to replace
    the within-session figures.
    \item \textbf{Word-level language modelling.} A dictionary- or trie-based
    decoder over the letter stream would resolve residual letter confusions at
    the word level, since context disambiguates letters that are individually
    ambiguous—mirroring how human readers parse fingerspelling.
    \item \textbf{Power hardening and battery operation.} Fixing the supply and
    wiring, adding decoupling, lowering the transmit power, and adding a
    listen-triggered sleep mode would enable untethered battery operation and
    reduce heat, allowing longer continuous use.
    \item \textbf{Larger and two-handed vocabulary.} Extending to a larger word
    set and, with a second instrumented hand, to two-handed signs would broaden
    practical usefulness.
    \item \textbf{Continuous recognition.} Moving from gated, segmented capture
    to continuous gesture spotting would make the device feel more natural,
    at the cost of a harder segmentation and modelling problem.
\end{itemize}

\section{Concluding Remarks}

The project delivers a complete, self-contained sign-language translation glove
and, just as importantly, a documented account of why it is built the way it
is. The system evolved through measured failures—antenna coupling, USB
throughput, BLE MTU, calibration drift, sensor resolution, and power
integrity—each of which was diagnosed and resolved with an understood
engineering trade-off rather than a guess. That process, and the
feature-representation conclusion it produced, is the work's main contribution
and the strongest foundation for the future work above.
